\scanpage{47}
\mspage{45}

\[
 X\longmapsto H_{M;x,y}(X)\overset{\mathrm{def}}{=}
 \mathrm{Hom}_{M(X)\Sigma_X^{-1}}(\deleted{\unsure{a}}\unsure{x|a}_X,\deleted{\unsure{b}}\unsure{y|b}_X),
\]

\deleted{\unsure{Non seulement}} \inserted{par la formule donc}

\[
 H(X)=\varinjlim_n
 \underbrace{\unsure{\Pi_0|\pi_0}\,\underline{\mathrm{Hom}}(X,K_n)}_{\substack{\Vert\\\mathrm{Hom}_{\mathrm{Cat\,hot}}(X,K_n)}} .
\]

Je ne sais donner d'interprétation de ce type ni de la catégorie $K_\infty$ elle-même, ni de son image dans Cat hot, ni de son image dans Hot.

Il est vrai que si je \unsure{restreint|restreins} le foncteur $H$ à \deleted{\unsure{ceux}} la sous-catégorie pleine de (Cat) formée des catégories \uline{finies} ($\mathrm{Ob}\,X$ et $\mathrm{Fl}\,X$ finis), alors on aura

\[
 \underline{\mathrm{Hom}}^{!}(X,K_\infty)\unsure{=|\simeq}\underline{\mathrm{Hom}}(X,K_\infty),
\]

et l'isom. précédent s'interprète comme

\[
 \underset{\substack{\Vert\,\mathrm{def}\\
 \mathrm{Hom}^{\unsure{\prime|\;}}_{M(X)\Sigma_X^{-1}}(\unsure{a|x}_X,\unsure{b|y}_Y)}}{H(X)}
 \simeq\Pi_0\,\underline{\mathrm{Hom}}(X,K_\infty)\overset{\text{déf.}}{\unsure{=|\simeq}}
 \mathrm{Hom}_{\mathrm{Cat\,hot}}(X,K_\infty).
\]

\deleted{\unsure{Ce}} Cela prouve du moins que la connaissance de $K_\infty$ est utile. Mais cette formule, pour $X$ restreint aux catégories \uline{finies}, ne détermine pas $K_\infty$ \deleted{``à homotopie} comme objet de Cat hot, à isom \inserted{unique} près dans cette catégorie --- \uuline{sauf} dans le cas \deleted{\unsure{bien}} \deleted{\unsure{douteux}} où $K_\infty$ est elle-$=$ \deleted{\unsure{don}} homotope à une \unsure{catégo}

% archive sheet number: 46
